% 本文件由 LaTeX 排版 Agent 根据有效 translation.json 自动生成。
% author=John Chrysostom; work_id=1911; title=论由屋顶缒下的瘫子讲道词

\chapter{\WorkTitle{论由屋顶缒下的瘫子讲道词}}

% structure: p0001 source=h1 target=omitted reason=page-title-already-rendered-by-parent

1. 不久前，我们偶遇了那位躺在池边床上的瘫子的事迹，发掘了一份丰硕而巨大的宝藏——不是因掘地而得，而是因潜入他的心肠而得：我们发现了一份宝藏，内含的不是金银宝石，而是坚忍、智慧、耐心以及对天主的极大望德，这比任何宝石或财富之源都更为宝贵。因为物质的财富易遭盗贼的计谋、诬告者的谗言、破门者的暴力、仆人的恶行所害，即便躲过了这一切，也常因激起嫉妒者的眼目而给拥有者带来最大的毁灭，从而生出无数麻烦的风暴。然而，属灵的财富却避开了所有这些祸害的机缘，超乎于一切此类辱骂之上，既嘲笑盗贼与破门者，也讥讽诽谤者、诬告者乃至死亡本身。因为死亡不能将它们与拥有者分离，相反，这拥有在死后反而更稳固地归于主人，它们伴随主人赴往另一世界，与他们一同移植于来世，成为离世者奇妙的辩护者，并使审判者对他们仁慈。\par

我们在这瘫子的灵魂中发现了这份财富的大量储备。你们是证人，曾以极大热忱汲取这份宝藏，却仍未将其耗尽。因为属灵财富的性质正是如此；它如同水泉，甚至比水泉更为丰盈，当有许多人汲取时，它便最为丰沛。因为当它进入任何人的灵魂时，它不被分割，也不减少，而是完整地临于每人，始终不被消耗，永不枯竭：这正是当时所发生的事。虽然已有如此多人接近这宝藏，且众人尽其所能地汲取——但我为何只提你们？因为从那时至今，它已使无数人变得富有，却仍保持其原初的完美。让我们不要因求助于这属灵财富之源而疲倦；而是尽可能现在也从中汲取，让我们凝视我们仁慈的主，凝视祂那忍耐的仆人。他患有不治之症已三十八年，长久受其折磨，却未怨声载道，未出亵渎之言，未控告造物主，而是勇敢且极其温顺地忍受了灾祸。这何以显明？你们说：因为圣经并未清楚告诉我们他先前的生活，只说他患病三十八年；并未加一词证明他没有表现不满、愤怒或急躁。然而，若有人仔细留意，非以好奇粗心之态审视，圣经也显明了这一点。因为当你们听见基督——对他而言是陌生人，且仅被视为常人——以如此大的温顺对他说话时，你们便能察觉他先前的智慧。因为当耶稣对他说：\Quote{你愿意痊愈吗？}他并未作出自然的回应：你看见我长久瘫卧病榻，竟问我是否愿意痊愈？你来是侮辱我的痛苦，责备我、嘲笑我、戏弄我的灾祸吗？他丝毫未说或想这类事，而是温顺地回答：\Quote{主，是的。}如今，若三十八年后他仍如此温顺柔和，而那时他所有理性的活力和力量都已衰颓，试想他在患难之初可能会怎样。须知，病患在疾病初期并不像经过长时间后那样难以取悦：当疾病绵延时，他们变得最倔强、最难忍受。但他多年后仍如此智慧，以如此克制的态度回答，显然在先前的时间里，他也以极大的感恩承受了那灾祸。\par

思及此，让我们效法同仆的忍耐：因为他的瘫痪足以振奋我们的灵魂：没有人会如此怠惰懒散，在观察了那灾祸的严重性后，不能勇敢忍受临于他的一切凶恶，即使它们比所有已知的苦难更难以忍受。因为不仅他的痊愈，连他的疾病也成了我们极大益处的缘由：他的治愈激励听众的灵魂赞美主，他的疾病与软弱鼓励你们忍耐，并敦促你们与他一样热心；或者更确切地说，它向你们显示了天主的慈爱。因为实际上，天主容许人患上如此疾病，且患病时间绵延如此之久，正是对他福祉极大关怀的标志。正如金匠将金子投入熔炉，任凭它在火中试炼，直至他看见它变得更纯净：同样，天主允许人的灵魂受患难考验，直至他们变得纯净透明，并从这筛炼过程中获得许多益处：因此，这是恩惠中最伟大的一种。\par

2. 那么，当试炼临于我们时，让我们不要烦乱，也不要沮丧。因为若金匠知道应将金子留在熔炉中多久，何时应取出，并不任其在火中毁坏烧尽：那么，天主更明白这事，当祂看见我们变得更纯净时，祂便从试炼中释放我们，以免我们因罪恶增多而倾覆跌倒。因此，当某种意外之事临于我们时，让我们不要抱怨或灰心；而是让那准确知道这些事的那位，按祂的意愿以火锻炼我们的心：因为祂为此行有益的目的，且着眼于受试炼者的利益。\par

因此，有一位智者劝诫我们说：\BibleQuote{我儿，你若前来事奉主，预备你的灵魂受诱惑，坚心忍耐，在患难时勿急躁；} \BibleRef{德 1:1} 他说\Quote{顺从祂}，\Quote{在一切事上}，因为祂知道何时适宜将我们从邪恶的熔炉中拔出。因此，我们应当处处顺从祂，常常感恩，乐意承受一切，无论祂赐予我们恩惠或惩戒，因为这也是一种恩惠。因为医生不仅当他给病人沐浴、滋养他、领他到宜人的花园时，而且当他使用烙和手术刀时，同样是医生；父亲不仅当他抚摸儿子时，而且当他逐出家门、斥责和鞭打他时，同样是父亲，并不亚于他称赞儿子的时候。因此，知道天主比一切医生更慈爱，不要过分好奇地追问祂的治疗，也不要向祂要求解释，但无论祂乐意释放我们，还是惩罚我们，让我们将自身同样献上；因为祂藉着两者都要领我们回归健康，并归向与祂的共融，祂知道我们各自的需要，什么对每个人有益，以及我们应当如何并怎样得救，祂就沿着那条道路引领我们。那么，让我们跟随祂所吩咐的任何地方，不要过分仔细地考虑祂命令我们走平坦容易的道路，还是艰难崎岖的道路；正如这瘫子的例子。久病的折磨炼净他的灵魂，使他被交付于苦难的火炼，如同投入某种熔炉，这确实是一种恩惠；但另一种恩惠丝毫不逊于此，即天主在试炼中与他同在，并赐予他极大的安慰。正是祂坚强了他，扶持了他，向他伸手，不容他跌倒。但当你听到是天主自己时，不要剥夺这瘫子应得的赞美，也不要剥夺任何受试炼而坚忍到底的人。因为即使我们无限智慧，即使我们比所有人更大更强，然而缺乏祂的恩宠，我们甚至连最寻常的诱惑也无法抵挡。为何我说到像我们这样微不足道卑劣的人？因为即使是一位\Person{保禄}、\Person{伯多禄}、\Person{雅各伯}或\Person{若望}，若他被剥夺了天主的帮助，也会轻易蒙羞、倾覆、俯伏在地。并且我要为这些人向你们宣读基督自己的话：因为祂对\Person{伯多禄}说：\BibleQuote{看，撒殚想要得到你们，好筛你们像筛麦子一样，但我已为你祈求，使你的信德不致丧失。} \BibleRef{路 22:31} \Quote{筛}的意思是什么？就是翻转、扭动、摇动、搅动、打破、折磨，如同簸扬时所发生的那样；但祂说，我已阻止了他，因为知道你们不能承受试炼，因为\Quote{使你的信德不致丧失}这话表明了若祂允许了，他的信德就会丧失。如果\Person{伯多禄}——他是如此炽爱基督，无数次为祂冒生命危险，跻身宗徒团体的最前列，且被他的师傅称为有福，并因此被称为\Person{伯多禄}，因为他坚持了坚定稳固的信德——如果基督允许魔鬼按他所愿地试探他，他就会被掳去，从信仰的承认中堕落，那么除却祂的帮助，还有谁能站立得住呢？因此\Person{保禄}也说：\BibleQuote{天主是忠信的，祂不许你们受试诱超过你们所能承受的，而且在试诱中也要开一条出路，使你们能忍受。} \BibleRef{格前 10:13} 因为祂不仅说不让超过我们力量的试炼加给我们，甚至在适合我们力量的试炼中，祂也临在扶持我们渡过，振作我们，只要我们首先贡献自己力所能及的方法，如热忱、对祂的盼望、感恩、坚忍、忍耐。因为不仅在与我们力量不相称的危险中，而且在相称的危险中，我们需要天主的帮助，才能勇敢站立；因为别处也记载：\Quote{正如基督的苦难多多加给我们，借基督我们的安慰也多多洋溢，使我们能以天主赐予我们的安慰，去安慰那些在患难中的人。} 那么，安慰这人的那位，就是允许试炼加给他的同一位。如今看祂在治愈后显示了何等的温柔。祂没有离开他走掉，而是在圣殿中找到他，说：\BibleQuote{看，你已经痊愈了，不要再犯罪，免得遭遇更糟的事。} \BibleRef{若 5:14} 因为若祂因恨他而允许惩罚，就不会释放他，也不会为他的将来安全作预备；但\Quote{免得遭遇更糟的事}这话是预先阻止将来的邪恶的言语。祂终止了疾病，但没有终止挣扎；祂驱除了软弱，但没有驱除对软弱的恐惧，以使已施行的恩惠保持稳固。这是慈心医生的做法，不仅终止当前的痛苦，而且为将来的安全预备，基督也这样做了，用对往事的回忆振作他的灵魂。因为看到当困扰我们的事离去时，对它们的回忆也常常随之离去，祂希望这回忆持续存留，便说：\Quote{不要再犯罪，免得遭遇更糟的事。}\par

3. 不仅如此，我们还能从祂看似责备的话语中看出祂的远见和体恤。因为祂并未公开揭露他的罪过，却告诉他，他所受的苦难是由于他的罪过；但祂并未指出是哪些罪过。祂没有说\Quote{你犯了罪}或\Quote{你违犯了诫命}，而是用一句简单的话指明了事实：\Quote{不要再犯罪了}；祂这样说了，足以提醒他，使他更警惕将来的事，同时向我们所有人彰显了他的忍耐、勇气和智慧，使他被迫公开哀悼自己的灾祸，并显明了祂对此人的关切：\Quote{因为当我来的时候，}他说，\Quote{别人先我下去了。}（见：\BibleRef{若 5:7}）然而祂并未公开揭露他的罪过。因为正如我们自己希望掩盖自己的罪过一样，天主更是如此：为此，祂在众人面前施行了治愈，却私下给予劝诫或忠告。因为祂从不公开显示我们的罪，除非祂看到人对罪麻木不仁。因为当祂说\Quote{我饿了，你们没有给我吃的；我渴了，你们没有给我喝的}（见：\BibleRef{玛 25:12}）时，祂现在这样说，是为了使我们将来不必听到这些话。祂在此世威胁、揭露我们，为的是在来世不必揭露我们；正如祂威胁要推翻尼尼微城（见：\BibleRef{纳 1:2}），正是为了不推翻它。因为如果祂想要公布我们的罪，祂就不会预先宣布祂要公布它们；但事实上，祂确实做了这个宣布，为的是使我们因惧怕暴露（即使不是因惧怕惩罚）而清醒，从而洁净自己全部罪过。同样的事也发生在圣洗圣事中：祂把人引到水泉，却不向任何人透露他的罪；然而祂公开地赐予恩惠，并使众人看见，而此人的罪除了天主自己和领受罪赦的人之外，无人知晓。瘫痪者的情况也是如此：祂在无见证人的情况下责备他，或者更确切地说，这话不仅是责备，也是一种辩护；祂仿佛在为长期如此折磨他而辩护，告诉他并向他证明，祂让他长久受苦并非无因无果，因为祂提醒了他的罪，并宣告了他软弱的原因。\Quote{因为找到他以后，}我们读到，\Quote{在圣殿里，祂对他说：\InnerQuote{不要再犯罪了，免得你遭遇更坏的事。}}\par

既然我们从以前那瘫痪者的记述中获得了如此多的益处，现在让我们转向玛窦福音所呈现给我们的另一位瘫痪者。因为在矿藏中，如果有人碰巧发现了一块金子，他会在同一地方进一步挖掘；我知道许多不经心阅读的人会以为四福音所呈现的是同一位瘫痪者，但事实并非如此；所以你们必须警醒，并密切注意此事。因为问题并非涉及寻常之事；这篇讲辞若得到恰当的解释，将有助于驳斥希腊人、犹太人以及许多异端者。因为所有这些人都会指责福音作者们互相争斗、有分歧；然而事实并非如此，上天不许！尽管外表不同，但运行在每位作者灵魂中的圣神的恩宠是同一的；哪里有圣神的恩宠，哪里就有爱、喜乐与和平；那里就没有战争、争论、纷争和冲突。那么我们如何能清楚表明这位瘫痪者不是另一位，而是不同的人呢？通过许多标记：地点、时间、季节、日子、治愈的方式、医生的来临以及被治愈者的孤独。这又怎样？有人会说：许多福音作者难道没有对其他征兆给出不同的叙述吗？是的，但提出不同的陈述是一回事，提出矛盾的陈述是另一回事；前者不会导致不和或争斗；但我们现在所面对的却是一个强烈的矛盾情况，除非能证明池边的瘫痪者与其他三位福音作者所描述的那位不是同一个人。为使你们明白不同的陈述与矛盾的陈述之间的区别，一位福音作者说基督背了十字架（见：\BibleRef{若 19:17}），另一位说基勒乃人西满背了它：但这并不导致矛盾或争斗。\Quote{你怎么说，}你说，\Quote{背和不背的陈述之间没有矛盾？}因为两者都发生了。当他们离开总督府时，基督正背着十字架；但前行时，西满从祂那里接过来背着。同样，关于那两个强盗，一位说两人都辱骂了，另一位说其中一人劝阻了那辱骂主的人（见：\BibleRef{路 23:40}）。然而这里又没有矛盾：因为这里同样两种情况都发生了；起初两人都行为恶劣，但后来当征兆出现，大地震动、岩石崩裂、太阳昏暗时，其中一人悔改了，变得更有节制，认出了被钉者并承认了他的王国。为了防止你认为这事是由某种内在驱使力量强迫发生的，并消除你的困惑，他向你展示那人在十字架上仍然保留着先前的邪恶，以便你看出他的悔改是由内而外、出于己心，在天主恩宠的帮助下完成的，从而他成为了一个更好的人。\par

4. 并且可以从福音中收集许多其他此类实例，它们似乎有矛盾的嫌疑，但实际上并无真正的矛盾；真相是，有些事件由这位作者记述，另一些由那位作者记述；或者，如果不是发生在同一时刻，一位作者叙述了较早的事件，另一位叙述了较晚的事件。但在目前的情况下，并无此类情形，而是我提到的众多证据向任何稍加留意此事的人证明，两次事件中的瘫子并非同一个人。这将是证明福音作者彼此和谐而非不一致的一个不轻的证据。因为若是同一个人，两段记述之间就有巨大分歧；但若是不同的人，所有争执的素材就被消除了。\par

那么现在请允许我陈述我断言这个人与那个人不是同一人的实际理由。理由是什么？一个在\Place{耶路撒冷}治愈，另一个在\Place{葛法翁}；一个在水池旁，另一个在某间房屋里：这是地点的证据；前者在庆节期间：这是特定节期的证据；前者患病三十八年：后者福音作者并未提及此类情况：这是时间的证据；前者在安息日治愈：这是日子的证据；因为如果此人也曾在安息日治愈，\Person{玛窦}不会沉默不语，在场的犹太人也不会保持沉默；因为那些因其他理由挑剔的人，即使一个人在安息日未得治愈，也会更猛烈地控告基督，如果他们从特殊日子的论点中得到了额外的把柄。此外，此人被带到基督面前；对另一个人，基督亲自来到他那里，而且没有人帮助他。他说：\Quote{主，我没有一个人。}而此人有许多人来帮助他，他们甚至把他从屋顶缒下去。祂在治愈另一个人的身体之前先治愈了他的灵魂：因为在治愈瘫痪之后，祂才说：\Quote{看，你已痊愈了，不要再犯罪。}但在此例中并非如此，而是在治愈了他的灵魂之后——因为祂对他说：\Quote{孩子，放心吧！你的罪赦了。}——然后才治愈了他的瘫痪。那么，此人与另一个人不是同一人，已通过这些证据清楚证明；但我们仍需转向叙述的开端，看看基督如何治愈这两人，为何方式不同：为何一个在安息日，另一个不在安息日；为何祂亲自来到一个人那里，却等待另一个人被带到祂面前；为何祂先治愈一个人的身体，另一个人先治愈灵魂。因为祂并不是漫无目的、随随便便地做这些事，因为祂是明智和审慎的。那么让我们留心观察祂施展治愈。因为对于医生，当他们使用手术刀或烧灼或对伤残瘫痪的病人进行其他操作、切除肢体时，许多人会围在病人和做这些事的医生周围；我们更应该在此事上如此行，因为医生更伟大，疾病更严重，是无法靠人力矫正、只能靠天主的恩宠医治的疾病。在之前的情况下，我们不得不看到皮肤被切开，脓液流出，血块涌动，忍受因景象而产生的极大不适，以及因那些遭受烧灼或切割的人所承受的痛苦而产生的极大悲伤；因为没有人如此铁石心肠，以至于能站在那些受苦的人旁边，听到他们的尖叫而不自身被打动和激动，不经历极大的精神压抑；但尽管如此，我们仍因渴望目睹手术而忍受这一切。但在此例中，无需看到任何此类景象，没有火的施用，没有器械的刺入，没有血液的流动，没有病人的痛苦或尖叫；其原因是医治者的智慧，祂不需要任何这些外在辅助，而是完全自足。因为只需祂发出一道命令，所有痛苦就止息了。令人惊奇的不仅是祂如此轻易地实现治愈，而且是无痛苦地，不给被治愈者带来任何麻烦。\par

既然这奇迹更大，这医治更重要，而给旁观者带来的喜悦丝毫不掺杂任何悲伤，那么就让我们现在仔细默想\Person{基督}在施行医治。祂上了船，渡到对岸，来到自己的城：看，有人把一个躺在床上的瘫子抬到祂跟前；\Person{耶稣}见了他们的信德，就对那瘫子说：\Quote{孩子！放心罢！你的罪赦了。}\BibleRef{玛 9:1}论信德，这些人不如百夫长，但比池旁的瘫子强。因为前者（百夫长）没有请医生，也没有把病人带到医生跟前，而是以天主的身份靠近祂，说：\Quote{只要祢说一句话，我的仆人就会痊愈。}\BibleRef{路 7:7}现在这些人没有请医生到家里，就此而言，他们与百夫长相等；但他们把病人带到医生跟前，就此而言，他们较为逊色，因为他们没有说：\Quote{祢只说一句话。}可是他们远比躺在池旁的那个人好。因为他说：\Quote{主，水动的时候，没有人把我放在池子里；}但这些人知道\Person{基督}既不需要水，也不需要池子，或任何那一类的东西。然而，\Person{基督}不仅释放了百夫长的仆人，也释放了另外两个人，使他们脱离疾病，且没有说：\Quote{因为你们提供的信德较小，你们所获得的医治也应与之相称；}而是让那表现更大信德的人带着赞美和荣誉离开，说：\Quote{我连在以色列也没有见过这样大的信德。}\BibleRef{路 7:9}对于比这人信德较小的人，祂没有给予称赞，却没有剥夺他的医治，不！甚至对那完全没有显示信德的人也是如此。但正如医生在治疗同一种疾病时，从一些人那里收取一百金币，从另一些人那里收取一半，从其他一些人那里收取较少，而从一些人那里分文不取：同样，\Person{基督}从百夫长那里领受了极大而不可言喻的信德，但从这个人那里领受较少，从那另一个人那里领受的甚至不到平常的量，然而祂医治了他们所有人。那么，祂为什么认为那个没有付出信德存款的人配得这恩惠呢？因为他没有显示信德，并非出于懒惰或灵魂麻木，而是由于不认识\Person{基督}，且从未听说过任何关于祂的奇迹，无论是小的还是大的。因此，这人获得了宽恕：事实上，圣史暗示了这一点，当他说：\Quote{因为他不知道那人是谁，}\BibleRef{若 5:13}但他第二次遇见祂时，才凭外貌认出了祂。\par

5. 的确，有些人说这个人之所以被治愈，仅仅是因为带他来的人相信了；但事实并非如此。因为\BibleQuote{祂看见他们的信德}不仅指那些带来的人，也指被带来的人。为何如此？你们说：\Quote{难道一个人因另一个人相信而被治愈吗？}依我看来，我不这样认为，除非由于年龄幼小或极度虚弱，他在某种程度上无法相信。那么，你们说，在客纳罕妇人的事例中，母亲相信了，女儿却被治愈了，这是怎么回事？百夫长相信了，他的仆人从病榻上起来并得以保全，这又是怎么回事？因为病人本身无法相信。请听客纳罕妇人说什么：\BibleQuote{我的女儿被魔鬼纠缠得好苦，有时跌进水里，有时跌进火里。}\BibleRef{玛 15:22} 一个心智被黑暗笼罩、被魔鬼附身、从未能自制、神志不清的人，怎能相信呢？同样，在百夫长的事例中，他的仆人卧病在家，不认识基督本人，也不知道祂是谁。那么他怎能相信一个素不相识、从未有过任何体验的人呢？但在我们眼前的事例中，我们不能这样说：因为瘫子相信了。这从哪里显明呢？从他接近基督的方式本身。不要只注意他们把人从屋顶缒下这件事；但要思想一个病人忍受这一切需要多大的毅力。因为你们一定知道，病人常是那样灰心丧气、难以取悦，以至于往往拒绝在病床上对他们施行的治疗，宁愿忍受疾病带来的痛苦，也不愿承受治疗引起的烦恼。但这个人有毅力走出屋外，被抬到集市中间，在人群面前展示自己。而且，病人通常宁愿死在疾病之下，也不愿暴露自己的个人不幸。但这个病人没有这样做，当他看到集会的地方挤满了人，通道被堵塞，避难所被阻挡时，他甘愿被从屋顶缒下。渴望是多么善于谋划，爱是多么富有创造力。\BibleQuote{因为凡寻找的，就找到；凡敲门的，就给他开门。}\BibleRef{路 11:10} 这个人没有对他的朋友们说：\Quote{这是什么意思？为什么这样吵闹？为什么推挤？我们等房子空了、集会散了再说；人群会散去，那时我们就可以私下接近祂，讨论这些事情。为什么你们要在众目睽睽之下暴露我的不幸，把我从屋顶上缒下，做出不合体统的行为？}这个人无论对自己还是对抬他的人都没有说这些，反而认为让这么多人见证他的治愈是一种光荣。不仅从这一情形可以辨认他的信德，而且从基督实际说的话也可以。因为在他被缒下并呈现在主前之后，基督对他说：\BibleQuote{孩子！放心吧，你的罪赦了。}他听了这些话，并没有愤慨，没有抱怨，没有对医生说：\Quote{你这是什么意思？我来是为了一件事得医治，你却医治另一件事。这是借口、托辞和无能的掩饰。你赦免看不见的罪？}他既没有说也没有想这些，而是等待着，允许医生采用祂所希望的治疗方法。正因如此，基督也没有到他那里去，而是等他前来，好向众人显明他的信德。难道祂不能使入口更容易吗？但祂没有这样做，为的是向所有人显明此人的热忱和炽热的信德。因为正如祂去到那病了三十八年的人那里，因为他没有人帮助，所以祂等待这个人到祂这里来，因为他有许多朋友，好通过这个人被带来显明他的信德，并通过去那里告诉我们另一个人的孤独，向所有人尤其是在场的人揭示这一个人的热忱和另一个人的忍耐。因为有些嫉妒和厌恶人类的犹太人习惯妒忌邻人所受的恩惠，并挑剔祂的奇迹，有时以特别的时节为由，说祂在安息日治病；有时以受惠者的生活为由，说：\BibleQuote{如果这个人是先知，祂必定知道摸祂的是怎样的女人}\BibleRef{路 7:39}，却不知道医生特有的标志是与软弱的人交往，并经常出现在病人身边，而不是避开他们或匆忙离开——这其实是祂对那些抱怨者明确说过的：\BibleQuote{健康的人不需要医生，而是有病的人需要。}\BibleRef{玛 9:12} 因此，为了防止他们再次提出同样的指控，祂首先证明那些到祂面前来的人因他们所表现的信德而配得治愈。为此，祂显露了一个人的孤独和另一个人的炽热信德与热忱；为此，祂在一个安息日治愈了那一个人，在另一个非安息日治愈了另一个人：以便当你们看到他们在另一天指控和责难基督时，你们会明白他们先前指控祂并非因为尊重法律，而是因为他们无法抑制自己的恶意。但为什么祂不首先着手治愈瘫子，反而说：\BibleQuote{孩子！放心吧，你的罪赦了}？祂这样做非常明智。因为医生惯于先消灭疾病的根源，然后再除去疾病本身。例如，当眼睛因某种恶性的体液和腐败的分泌物而痛苦时，医生放弃对视力的任何治疗，转而注意头部，那里是软弱根源和起源所在：基督正是这样做的：祂首先压制了邪恶的源头。因为一切邪恶的源头、根源和母亲是罪的性质。正是它削弱了我们的身体；正是它带来了疾病；所以此时祂说：\BibleQuote{孩子！放心吧，你的罪赦了。}而在另一处祂说：\BibleQuote{看，你已经痊愈了，不要再犯罪，免得遭遇更糟的事}，表明这两种疾病都是罪的产物。在起初和开端，疾病作为罪的后果袭击了\Person{加音}的身体。因为他杀害兄弟之后，在那邪恶行为之后，他的身体患了瘫痪。因为颤抖与瘫痪是同一回事。当支配生物的力量变得虚弱，不能再支撑所有肢体时，便剥夺了它们天然的控制力，然后肢体松弛，颤抖并眩晕。\par

6. 保禄也证明了这一点：因为当他责备格林多人某种罪时，他说：\BibleQuote{因此你们中间有许多软弱和患病的。}因此基督首先除去了邪恶的根源，说了\BibleQuote{孩子！放心吧，你的罪赦了，}他振奋精神，唤醒沮丧的灵魂：因为言语成为有效的原因，进入良心，抓住了灵魂本身，从中驱逐了一切忧愁。因为没有什么比免于自责更能创造快乐并提供信心。因为哪里罪得赦免，哪里就有子女的身份。至少我们也是如此，直到在圣水之池中洗去我们的罪，我们才能称天主为父。正是当我们从那里上来，脱去那罪恶的重担时，我们说\BibleQuote{我们的天父，你在天上。}但在那个患病三十八年的人身上，他为什么没有这样做，而是首先医治了他的身体？因为在那长时间里，他的罪已经耗尽：因为考验的严重性能减轻罪的担子；正如我们读到拉匝禄的情况，他完全承受了恶事，随后得到了安慰；又在另一处我们读到：\BibleQuote{你们安慰我的百姓，向耶路撒冷的心说话，因她从主的手中领受了双倍的罪罚。}\BibleRef{依 40:1} 先知又说：\BibleQuote{上主赐给我们平安，因你已向我们偿还了一切，}\BibleRef{依 26:12} 这表明惩罚和刑罚产生罪的赦免；我们可以从许多段落证明这一点。在我看来，他之所以对那人只字不提罪的赦免，只确保他在将来安全，是因为他的罪已通过长期的疾病得到了惩罚：如果不是这个原因，那就是因为他尚未达到对基督的高度信仰，以致主首先针对较小的、明显可见的需要，即身体的健康；但在另一个人身上，他没有这样做，而是因为这个人有更大的信德和更高尚的灵魂，他首先针对更危险的疾病说话：此外还为了显示他与圣父同等的地位。正如在前一个事例中，他在安息日治病，因为他想引导人们离开犹太人守安息日的方式，并利用他们的责难来证明自己与圣父同等：同样，在这个事例中，他预先知道他们要说什么，就说出这些话，以便利用它们作为起点和借口，证明他与圣父同等的地位。因为当没有人提出控告或指控时，主动开始谈论这些事是一回事；而当其他人提供机会时，以辩护的形式和顺序来进行同样的工作则是另一回事。因为前一种证明的性质对听者来说是绊脚石；但后一种较不令人反感，更容易被接受，我们到处看到他这样做，并且显示他的同等不凭言语而凭行为。无论如何，这正是圣史所暗示的，当他说犹太人迫害耶稣，不仅因为他违反了安息日，而且因为他说天主是他的圣父，使自己与天主同等，\BibleRef{若 5:16} 这是更大的事，因为他通过行为的证明实现了这一点。那么，那些嫉妒和邪恶的人，以及那些处处寻找把柄的人，是如何做的呢？他们说：\BibleQuote{这人为什么说亵渎的话？}因为\BibleQuote{除了天主独一以外，无人能赦罪。}\BibleRef{谷 2:7} 正如他们因他违反安息日而在那里迫害他，并利用他们的责难以辩护的形式宣告他与圣父同等，说\BibleQuote{我的圣父直到现在工作，我也工作，}\BibleRef{若 5:17} 同样，在这里，他从他们的指控出发，以此证明他与圣父完全相似。他们说了什么？\BibleQuote{除了天主独一以外，无人能赦罪。}既然他们自己定下了这个定义，他们自己引入了规则，他们自己宣布了律法，他就继续用他们自己的话来缠住他们。他说：\Quote{你们承认了，罪的赦免是天主独有的属性：因此，我的同等是无可置疑的。}不仅这些人如此宣告，先知也这样说：\BibleQuote{谁神相似你？}然后，指示他的特有属性，他接着说\BibleQuote{赦免罪过，宽恕不义。}\BibleRef{米 7:18} 如果其他人显现做同样的事，他也是天主，正如那一位是天主。但让我们观察基督如何与他们争辩，多么温良、柔和，充满温柔。\BibleQuote{看，有些经师心里说：这人亵渎了。}他们没有说出这句话，没有用舌头宣扬，而是在心的隐秘处思量。基督如何做？他在与医治瘫子身体有关的证明之前，公开了他们的秘密思想，想向他们证明他天主性的能力。因为显示人心中的秘密是天主独有的属性，是他天主性的记号，圣经说：\BibleQuote{惟独你知晓人心。}\BibleRef{列上 8:39} 你看到这个\Quote{独一}一词并非为了对比圣子与圣父。因为如果圣父独知人心，圣子如何知道思想中的秘密？经上说：\Quote{他自己知道人里面有什么；}\BibleRef{若 2:25} 保禄在证明知晓隐秘之事是天主特有的属性时说：\Quote{那洞察人心者，}\BibleRef{罗 8:27} 表明这一表达等同于\Quote{天主}的称号。因为正如当我说\Quote{那降雨的说}，我通过提及作为来指明天主，别无他指，因为这作为只属于他独一；当我说\Quote{那使太阳升起的}，虽不加上\Quote{天主}一词，我仍通过提及作为来指明他；同样，当保禄说\Quote{那洞察人心的}，他证明了洞察人心是天主独有的属性。因为若这一表达在指明所指者时不具有与\Quote{天主}之名同等的力量，他就不会绝对地、单独地使用它。因为若这权能是他与某受造物共同分享的，我们就不会知道所指为谁，权能的共有造成听者心中的混乱。如此，既然这似乎是圣父特有的属性，却也在圣子身上彰显，由此圣子的同等成为不容置疑，因此我们读到：\BibleQuote{你们为什么心里怀着恶念？什么更容易呢？是说：你的罪赦了，还是说：起来行走？}\par

7. 再者，祂又提供了第二个证据，证明祂有权赦罪。因为赦罪远比医治身体更为伟大：灵魂比身体伟大多少，赦罪就比医病伟大多少。瘫痪是身体的疾病，同样，罪是灵魂的疾病。但罪虽然更大，却不可见；而病虽然较小，却是明显的。既然祂要用那较小的事来证明更大的事，表明祂是因他们的软弱而行这事，并迁就他们脆弱的状况，祂就说：\Quote{哪一样更容易呢？是说\InnerQuote{你的罪赦了}，还是说\InnerQuote{起来行走}呢？}那么，祂为何要为他们而从事那较小的行动呢？因为明显的事能更清晰地提供证据。因此，祂没有先让那人起来，而是先对他们说：\Quote{但为叫你们知道，人子在地上有权赦罪——就对瘫痪者说——起来，行走吧！}好像祂在说：赦罪固然是更大的标记，但为了你们，我也加上这较小的，因为这在你们看来似乎是另一个的证明。正如在另一场合，祂称赞百夫长说：\Quote{只要你说一句话，我的仆人就会痊愈：因为我也对这个人说\InnerQuote{去}，他就去；对另一个人说\InnerQuote{来}，他就来。}祂用所发出的赞扬来肯定百夫长的看法。又如，当祂责备犹太人在安息日挑剔祂，说祂违法时，祂证明自己有权改变法律。同样，在此处，当有人说\Quote{祂应许只属于圣父的事，使自己与天主平等}时，祂责备并指控他们，用行动证明祂并未亵渎，从而为我们提供了无可辩驳的证据，表明祂能与生养祂的圣父做同样的事。至少请注意祂乐意建立这一事实的方式：即只属于圣父的事，也属于祂自己。因为祂不仅让瘫痪者起来，而且说：\Quote{但为叫你们知道，人子在地上有权赦罪。}因此，祂的奋斗和热切愿望，就是首先证明祂拥有与圣父相同的权柄。\par

8. 因此，让我们谨慎地持守这一切——无论是昨天还是前天所说的——并恳求天主使它们稳固地存留在我们心中。同时，我们也当贡献自己的热忱，恒常在此聚集。因为这样，我们就能保全先前所讲的真理，并增添新的储备；即使其中有些因时间久远而从记忆中滑落，我们也能借助不断的教导轻易恢复。如此，不仅教义能保持纯正无缺，我们的生活也将得到许多细心的照料，并能愉快而喜乐地度过此生。因为无论何种痛苦压迫我们的灵魂，当我们来到这里，就能轻易摆脱；看，如今基督也在此，凡以信德亲近祂的人，必能轻易从祂那里获得医治。假如有人被恒久的贫困所困扰，缺乏必需的食物，常常饿着肚子上床，若他来到这里，听到保祿说他忍饥挨饿、赤身露体，而且这不仅是一两天或三天，而是经常如此（至少他在说\Quote{直到此时此刻，我们又饥又渴，又赤身露体}时暗示了这一点）\BibleRef{格前 4:11}，他将获得极大的安慰，从而明白：天主并未因憎恨或遗弃他而允许他陷于贫困；因为若这是憎恨的结果，祂就不会允许这事临到保祿——这位尤其蒙祂挚爱的人身上。但祂是出于温柔的爱和眷顾，并以引导他达到更高属灵智慧的方式允许了这事。另一个人，身体被疾病和无数痛苦缠绕？这些瘫痪者的状况，加上那蒙福而勇敢的保祿的门徒——他不断受疾病折磨，从未从长期的虚弱中得到喘息——甚至如保祿所说：\Quote{为了你的胃和你的屡次虚弱，稍微用点酒吧。}\BibleRef{弟前 5:23}这里他不仅说虚弱本身。另一个人，受诬告而在公众中名声败坏，这事不断地烦扰并噬咬他的灵魂；他进入此地，听到：\Quote{当人们辱骂你们，捏造各样坏话毁谤你们时，你们是有福的：你们欢喜踊跃吧，因为你们在天上的赏报是丰厚的。}\BibleRef{玛 5:11}然后他将抛开一切沮丧，领受各种快乐；因为经上记载：\Quote{当人们因你们的名字为邪恶而弃绝你们时，你们要踊跃欢喜。}\BibleRef{路 6:22}同样，天主这样安慰受毁谤的人，又用另一种方式警告毁谤者说：\Quote{人所说的每句恶言，无论是好是坏，他们都要为此交账。}\BibleRef{玛 12:36}\par

另有一人也许失去了一个小女儿或儿子，或他的一位亲属；他也来到这里，聆听\Person{保禄}对今世生活的呻吟，渴望那将来的，因在这世界的旅居而受压抑；当他听见\Person{保禄}说：\BibleQuote{论到长眠的人，我不愿意你们不知道，免得你们悲伤，像那些没有望德的人一样。}\BibleRef{得前 4:13} 他没有说\Quote{论到死亡的人}，却说\BibleQuote{论到长眠的人}，证明死亡就是睡眠。正如我们看见有人睡觉时，不会惊慌或忧愁，因为他必定会起来；同样，当我们看见有人死了，也不要惊慌或沮丧，因为这也是睡眠，虽是较长的睡眠，却仍是睡眠。祂用\Quote{长眠}这个称呼，安慰了哀悼者，也驳斥了不信者的责难。你若对已逝者过度哀伤，就像那不信者一样，没有复活的望德。他哀伤固然有理，因为他不能对将来之事运用属灵智慧；但你既已领受了关于未来生命如此强有力的证据，为何与他同陷于这同样的软弱？因此经上写道：\BibleQuote{论到长眠的人，我们不愿意你们不知道，免得你们悲伤，像那些没有望德的人一样。}\par

不仅从新约，就是从旧约也能获得丰富的安慰。当你听到\Person{约伯}在失去财产、失去牲畜、失去不是一两个或三个，而是一整队正值盛年的儿子，以及他展现的伟大灵魂之后，即使你是最软弱的人，也能轻易悔改并重振勇气。因为，人啊，你常守护你生病的儿子，看见他躺在床上，听见他临终的话语，在他断气时站在身旁，亲手合上他的眼睛，闭上他的口；但\Person{约伯}并未目睹儿子们垂死的挣扎，没有看见他们咽下最后一口气，房屋成了他们共同的坟墓，在同一张桌上，脑浆与鲜血倾洒，木片、瓦砾、尘土与碎肉混杂在一起。然而，经历了如此巨大的灾祸，他并没有怨天尤人，而是说什么呢？——\BibleQuote{主赏赐的，主收回；主看着好，就如此成就了；愿主的名永受赞美。} 让我们在每一件临到我们的事上也说出这话，无论是损失财产、身体病弱、受辱、被诬告，或任何其他人类遭遇的祸患，我们都当说：\BibleQuote{主赏赐的，主收回；主看着好，就如此成就了；愿主的名永受赞美。} 如果我们修习这属灵的智慧，即使承受无数苦难，也绝不会经历任何祸患；因为得益必将大于损失，善必将胜于恶；用这些话，你将使天主怜悯你，并保卫自己免受\Person{撒殚}的暴虐。因为一旦你的舌头说出这些话，\Person{撒殚}立刻就远离你；\Person{撒殚}离开，沮丧的阴霾也消散，折磨我们的思绪也随之逃去，与他一同匆忙离去；除此之外，你还会在此世和天堂赢得一切祝福。你有令人信服的先例，就是\Person{约伯}和宗徒的事迹：他们为了天主的缘故轻视今世的烦恼，获得了永远的福乐。因此，让我们充满信心，在所有临到我们的事上欢喜感恩于仁慈的天主，使我们能宁静地度过今世，并获得将来的福乐，赖我们的主耶稣基督的恩宠和慈爱，愿光荣、尊威和权能归于祂，从今世直到永远。阿们。\par

\SourceRecord{Translated by W.R.W. Stephens.；Nicene and Post-Nicene Fathers, First Series；Vol. 9.；Edited by Philip Schaff.；Buffalo, NY: Christian Literature Publishing Co.,；1889.；<http://www.newadvent.org/fathers/1911.htm>.}
